% Clause 5: source PDF pages 32--33 / printed pages 20--21.
% Fidelity pass 03: source-calibrated subclause spacing and paragraph rhythm.
% The source uses an additional blank 6 pt between a rule statement and its
% explanatory paragraph (and between the three Rule 3 explanatory paragraphs).
% Forced line breaks preserve source line membership under Tinos.

\newcommand{\federationrulesubclause}[1]{%
    \par\vspace{11.4pt}%
    {\sffamily\bfseries\fontsize{10.98pt}{13.2pt}\selectfont #1\par}%
    \vspace{6.6pt}%
}

\vspace{2.1pt}
\clause{5. Federation rules}

This clause describes the five rules that apply to HLA federations. Rationale for the inclusion of each rule is\linebreak
located in Annex A.

\federationrulesubclause{5.1 Rule 1}

\textit{Federations shall have an HLA FOM, documented in accordance with the HLA OMT.}
\vspace{6pt}

All data to be exchanged in accordance with the HLA shall be documented in a FOM. The FOM is specified\linebreak
using one or more FOM modules and one MOM and Initialization Module (MIM). A FOM shall document\linebreak
the agreement among federates on data to be exchanged using the HLA services during federation execution\linebreak
and the minimal set of conditions of the data exchange (e.g., updates to be sent when changes exceed a\linebreak
certain value). As such, a FOM is an essential element in defining a federation. The HLA does not prescribe\linebreak
which data are included in a FOM (this task is the responsibility of the federation user and developer).\linebreak
FOMs shall be documented in the format prescribed in IEEE Std 1516.2-2010 to support reuse of a FOM by\linebreak
new users for their purposes.

\federationrulesubclause{5.2 Rule 2}

\textit{In a federation, all simulation-associated object instance representation shall be in the federates, not in the\linebreak
RTI.}
\vspace{6pt}

In the HLA, the responsibility for maintaining the values of HLA object instance attributes shall take place\linebreak
in the joined federate. In an HLA federation, all joined federate-associated instance attributes shall be owned

\newpage

by federates, not by the RTI. However, the RTI may own instance attributes associated with the federation\linebreak
MOM. The RTI may use data about instance attributes and interactions to support RTI services (e.g.,\linebreak
Declaration Management), but these data are merely used by the RTI, not changed.

\federationrulesubclause{5.3 Rule 3}

\textit{During a federation execution, all exchange of FOM data among joined federates shall occur via the RTI.}
\vspace{6pt}

The HLA federate interface specification (see IEEE Std 1516.1-2010) specifies a set of interfaces to services\linebreak
in the RTI to support coordinated exchange of instance attribute values and interactions in accordance with a\linebreak
federation's FOM. Under the HLA, intercommunication of FOM data among joined federates participating\linebreak
in a given federation execution shall be executed by the exchange of data via the RTI services. Based on the\linebreak
FOM, joined federates shall identify to the RTI what information they will provide and require, along with\linebreak
instance attribute and interaction data corresponding to the changing state of object instances in the joined\linebreak
federate. The RTI shall then provide the coordination, synchronization, and data exchange among the joined\linebreak
federates to permit a coherent execution of the federation.
\vspace{6pt}

Ensuring that the right data are provided at the right times and that the data are used in a substantively\linebreak
correct way shall be the responsibility of the joined federates. The RTI shall ensure that the data are\linebreak
delivered using joined federates in accordance with their declared requirements (e.g., which data, reliability\linebreak
of transport, event ordering) to provide a common view of shared data across a given federation execution,\linebreak
as specified in the FOM.
\vspace{6pt}

RTI services shall be used to ensure that the coordination needs of the distributed applications (federations)\linebreak
are met in a consistent way across all participants in a federation and over the life of a federation execution.

\federationrulesubclause{5.4 Rule 4}

\textit{During a federation execution, joined federates shall interact with the RTI in accordance with the HLA\linebreak
interface specification.}
\vspace{6pt}

The HLA provides a specification for a standard interface between the federate application and the RTI.\linebreak
Joined federates shall use this standard interface to access RTI services (see IEEE Std 1516.1-2010). The\linebreak
specification shall define how federate applications interact with the infrastructure. However, because the\linebreak
interface and the RTI will be used for a wide variety of applications requiring data exchange of diverse\linebreak
characteristics, the interface specification says nothing about the specific federate data to be exchanged over\linebreak
the interface.

\federationrulesubclause{5.5 Rule 5}

\textit{During a federation execution, an instance attribute shall be owned by, at most, one joined federate at any\linebreak
given time.}
\vspace{6pt}

The HLA allows for different joined federates to own different attributes of the same object instance (e.g., a\linebreak
simulation of an aircraft might own the location of the airborne sensor, whereas a sensor system model\linebreak
might own other instance attributes of the sensor). To ensure data coherency across the federation, at most,\linebreak
one joined federate may own any given instance attribute of an object instance at any given time. Joined\linebreak
federates may request that the ownership of instance attributes be acquired or divested, dynamically, during\linebreak
federation execution. Thus, ownership can be transferred, dynamically during execution, from one joined\linebreak
federate to another.
